\section{Métodos matemáticos e estratégias de controle}

A literatura relacionada às operações de rendezvous, berthing e robótica espacial emprega diferentes métodos matemáticos e estratégias de controle para representar o movimento relativo, descrever a dinâmica do sistema e conduzir manobras orbitais e tarefas de manipulação. Essa diversidade reflete o caráter multidisciplinar do problema, que envolve dinâmica orbital relativa, atitude, modelagem de manipuladores, controle, otimização e, em trabalhos mais recentes, técnicas computacionais orientadas por dados \cite{fonseca2016dynamics,arantes2010far,fiori2024modeling,santos2022machine,santos2023physics}.

Entre os métodos matemáticos mais recorrentes, destaca-se a formulação lagrangiana para a modelagem de sistemas acoplados. No trabalho de da Fonseca, Unfried e Lima, essa abordagem é empregada com multiplicadores de Lagrange e função dissipativa de Rayleigh para descrever a dinâmica de uma espaçonave do tipo manipulador durante a aproximação de curta distância, incluindo forças e torques de reação nas juntas e seus efeitos sobre a plataforma \cite{fonseca2016dynamics}. Essa formulação é particularmente relevante porque permite representar, de forma consistente, os acoplamentos entre os corpos do sistema e os efeitos dinâmicos decorrentes do movimento do manipulador \cite{fonseca2016dynamics}.

No tratamento do movimento relativo orbital, as equações de Hill--Clohessy--Wiltshire continuam desempenhando papel importante como ferramenta matemática para descrever a dinâmica linear relativa entre perseguidor e alvo em órbitas circulares ou quase circulares. No trabalho de \cite{arantes2010far}, essas equações compõem a base da análise das manobras far e proximity de satélites de serviço, em conjunto com estratégias de transferência orbital e aproximação ao satélite cliente \cite{arantes2010far}. De modo semelhante, o estudo de da Fonseca, Unfried e Lima também utiliza as equações de Hill como parte da formulação empregada nas simulações da aproximação final \cite{fonseca2016dynamics}.

Do ponto de vista do controle, observam-se tanto estratégias clássicas quanto formulações geométricas mais recentes. No estudo de da Fonseca, Unfried e Lima, técnicas LQR e PID são empregadas, respectivamente, para o movimento orbital linear relativo e para a estabilização da atitude durante a operação do manipulador \cite{fonseca2016dynamics}. Já no trabalho de Fiori et al., o rendezvous automatizado sob restrições posicionais é tratado com potenciais virtuais e um formalismo baseado em grupos de Lie, buscando lidar com obstáculos e com a orientação necessária à aproximação final \cite{fiori2024modeling}. Essas diferenças mostram que a escolha da estratégia de controle depende do tipo de missão, do modelo adotado e das restrições operacionais consideradas \cite{fiori2024modeling,fonseca2016dynamics}.

Além dessas abordagens, métodos de otimização e técnicas de aprendizado de máquina também têm sido incorporados como ferramentas auxiliares. Os trabalhos de \cite{santos2022machine,santos2023physics,santos2022uncertainty} indicam que essas técnicas podem apoiar o planejamento de trajetória e a resolução de problemas inversos, melhorando estimativas iniciais, convergência e eficiência computacional. Ainda assim, tais abordagens aparecem, em geral, como complemento às formulações físicas do sistema, e não como substituição da modelagem dinâmica e cinemática clássica \cite{santos2022machine,santos2023physics,santos2022uncertainty}.

De modo geral, observa-se que os métodos matemáticos e as estratégias de controle empregados na literatura abrangem desde formulações energéticas e equações lineares de movimento relativo até técnicas clássicas de controle, potenciais virtuais, grupos de Lie, otimização e aprendizado de máquina. Para a presente dissertação, essa revisão é relevante porque evidencia os principais instrumentos teóricos e computacionais disponíveis para sustentar a modelagem integrada da dinâmica relativa, da atitude da espaçonave perseguidora e do comportamento do manipulador robótico em operações de berthing.